\section{Infinity as a change of rules}\label{sec:infinity}

Theorem~\ref{thm:hfsat} left one axiom out of the free column: Infinity. This section shows that
the omission is forced. No fixed finite packaging rule, however often it is iterated from
hereditarily finite sets, produces an infinite set. Infinity is therefore the first place where
the construction must be told to continue: it is a change of rule, not a longer use of the old
one.

\subsection{Finite packaging rules}

The result is informative only if the class of excluded rules is broad enough to contain the
operations of \S\ref{sec:free}. Those operations draw new members not just from the current
family $A$ but also from inside its elements (the union and gate clauses). The right class is
therefore the following.

\begin{definition}[Finitely generated completion rule]\label{def:fg}
An operator $G$ on families of sets is a \emph{completion rule} if $A\subseteq G(A)$ for every
$A$. It is \emph{finitely generated} if every $x\in G(A)\setminus A$ is a \emph{finite} set with
\[
  x\;\subseteq\;A\cup\bigcup_{a\in A}\TC(a).
\]
\end{definition}

\begin{lemma}[The packaging operator is finitely generated; \inhouse]\label{lem:Ffg}
Extend $F$ (Definition~\ref{def:F}) to arbitrary families by requiring the union, gate and
transport outputs to be finite. The extended operator is finitely generated and agrees with $F$
on every $A\subseteq\HF$, where the finiteness requirement holds automatically.
\end{lemma}

The finiteness requirement is essential. Without it, the extension could form the union of an
infinite set, or one of its infinite subsets, and nothing would confine it to $\HF$.

\begin{lemma}[$\HF$-stability; \inhouse]\label{lem:hfstable}
If $G$ is finitely generated and $A\subseteq\HF$, then $G(A)\subseteq\HF$. Hence every finite
iterate $G^n(A_0)$ of a seed $A_0\subseteq\HF$ stays inside $\HF$.
\end{lemma}

\begin{proof}
A new element $x$ is a finite set whose members lie in $A$ or in the transitive closure of a
member of $A$. The transitive closure of a set in $\HF$ consists of sets in $\HF$, so every member
of $x$ is in $\HF$, and so is $x$. Induct on $n$.
\end{proof}

\subsection{Non-attainment}

\begin{theorem}[No finitely generated iteration reaches an inductive set;
\inhouse]\label{thm:nonattain}
Let $G$ be a finitely generated completion rule and $A_0\subseteq\HF$. Put $A_{n+1}=G(A_n)$ and
$A_\omega=\bigcup_n A_n$. Then:
\begin{enumerate}\itemsep0pt
  \item no $A_n$ contains an inductive set as an element;
  \item $A_\omega\subseteq\HF$, so $A_\omega$ does not contain $\omega$ or any other inductive
  set either; admitting $\omega$ requires a step outside the finitely generated class;
  \item if $G$ is moreover idempotent, the iteration stabilizes after one step:
  $A_n=A_1$ for all $n\ge1$.
\end{enumerate}
\end{theorem}

\begin{proof}
By Lemma~\ref{lem:hfstable} every $A_n$ is a subfamily of $\HF$, and no element of $\HF$ is
inductive (Theorem~\ref{thm:hfsat}). This gives (1), and (2) follows because a union of
subfamilies of $\HF$ is a subfamily of $\HF$; a step that added $\omega$ would add an infinite
set, which Definition~\ref{def:fg} forbids. (3) is immediate from $G(G(A_0))=G(A_0)$.
\end{proof}

Theorem~\ref{thm:nonattain} is a statement about objects produced by rules. It does not say that
finite axiom systems, or formal theories in general, cannot prove that an infinite set exists;
that is a different question, about syntax rather than construction.

\subsection{Infinity is a package change}

\begin{corollary}[The first purchase; \inhouse]\label{cor:packagechange}
Passing from $\HF$ to a family that contains an inductive set cannot be achieved by iterating any
finitely generated completion rule from any seed inside $\HF$. In the framework's terms, admitting
$\omega$ is a \emph{package change}: the new object is not produced by running the fixed package
longer.
\end{corollary}

We call the corresponding commitment \textbf{C-INF}: the admission of a completed infinite
collection, in the package-change sense of \citet{SBTIncompleteness}. We do not assert, derive
or evaluate \textbf{C-INF} here. We have shown only that it cannot be avoided by iteration.

\begin{remark}[Saturation and incompleteness; \metag]\label{rem:saturation}
Corollary~\ref{cor:packagechange} is one instance of a general pattern: a fixed completion rule
that adds only finite objects may keep adding them forever, but by iterating itself it cannot
produce an object of a new kind, here an infinite set. There
is a deliberate resonance with incompleteness. G\"odel's theorems \citep{Godel1931} say that a
consistent, effectively axiomatized theory containing enough arithmetic does not prove every
arithmetical truth, and (under the usual provability conditions) does not prove its own
consistency. Theorem~\ref{thm:nonattain} says that a fixed finite packaging rule does not reach
the infinite. Both describe a closed apparatus, run to saturation, with something deliberately out
of reach. The other half of the picture---what adjoining new content does---is the subject of
\S\ref{sec:forcing}. The resemblance between the two halves is a reading, not a metatheorem.
\end{remark}
